\subsubsection{A common six-unit heat slab: all-order division-free comparison}
\label{subsec:sixbanddivisionfree}
The preceding all-order heat positivity band can be made
\emph{explicit} without computing a separate Laguerre
threshold for each order. More importantly, the same
calculation bounds the coupled parabolic source directly,
without dividing by a Laguerre expression at a point
where it might be small or vanish.

For $0\le t\le1/2$ write
$\mu_j(t)=\int_0^\infty u^j e^{tu^2}\Phi(u)\,du$.
Let $A_n(t)=L_{n,t}(0)>0$ for the half-line transform
$H_t$ and $k_n=(n+1)(2n+1)$.

\begin{theorem}[Explicit common division-free heat band]
\label{thm:sixband}
For the genuine Jacobi-theta heat family,
every $n\ge1$, every $t\in[0,1/2]$ and
every $x\in\mathbb R$ satisfy
\begin{equation}\label{eq:sixbandlower}
\boxed{
L_{n,t}(x)\ge
A_n(t)\left(1-\frac{x^2}{48}\right).}
\end{equation}
Define the positive coefficient
\[
c_n(t)=4k_n\,\frac{A_{n+1}(t)}{A_n(t)}.
\]
Then for \emph{all} real $x$
\begin{equation}\label{eq:sixbandresidual}
\boxed{\begin{aligned}
k_nL_{n+1,t}(x)-c_n(t)L_{n,t}(x)
&\le k_nA_{n+1}(t)
\left(\frac{x^2}{12}-3\right).
\end{aligned}}
\end{equation}
In particular the required source inequality
\begin{equation}\label{eq:sixbandsource}
\boxed{k_nL_{n+1,t}(x)\le
c_n(t)L_{n,t}(x)}
\end{equation}
holds simultaneously for \emph{every order and heat time}
whenever $|x|\le6$.
The coefficient involves only strictly positive
central correlation moments. No division by the
potentially vanishing noncentral $L_{n,t}(x)$
is used anywhere.

\end{theorem}
\begin{proof}
The earlier validated two-mode moment certificate
(Theorem~\ref{thm:firstcompactrectangle},
equation~\eqref{eq:rectangle-moment-bounds}) gives
\[
\mu_{0,2}(0)>\frac1{20},\qquad
\mu_{2,2}(1/2)<\frac1{1000}.
\]
The omitted $m\ge3$ theta terms are positive.
The uniform heat-tail estimate
\eqref{eq:uniformheattail} with $k=2$,
applied to the half-line transform at the origin,
bounds their second moment by
\[
0\le\mu_2(t)-\mu_{2,2}(t)
\le\frac3{16}q_3e^{-q_3+1/16},
\quad q_3=9\pi.
\]
The elementary inequalities
$27<q_3<198/7$, $q_3-1/16>24$
and $e^3>20$ give the exact rational estimate
\[
\frac3{16}q_3e^{-q_3+1/16}
<\frac3{16}\frac{198/7}{20^8}
<10^{-8}.
\]
As $\mu_j(t)$ is increasing in $t$
for each even $j$, we conclude
\begin{equation}\label{eq:sixmomentgate}
\boxed{
\frac{\mu_2(t)}{\mu_0(t)}
<\frac{1/1000+10^{-8}}{1/20}
=\frac{100001}{5000000}<\frac1{48}.}
\end{equation}

The heat tilt changes the potential
$-\log\Phi$ only by $-tu^2$,
so Newman's radially monotone
curvature criterion applies at every $t$
\cite{NewmanGHS}. The reverse-likelihood-ratio
argument of Proposition~\ref{prop:radial}
therefore gives
\[
\frac{B_{n,t}}{A_{n,t}}
\le\frac{B_{0,t}}{A_{0,t}}
=2\frac{\mu_2(t)}{\mu_0(t)}<\frac1{24}.
\]
Here $A_{n,t},B_{n,t}$ are the
positive correlation moments of the
half-line transform, differing from
those for the full-line transform by a
common factor of $1/4$.
Using $\cos(xw)\ge1-x^2w^2/2$ yields
\[
L_{n,t}(x)\ge A_n(t)
-\frac{x^2}{2}B_{n,t}
\ge A_n(t)\left(1-\frac{x^2}{48}\right),
\]
proving \eqref{eq:sixbandlower}.
Also
$L_{n+1,t}(x)\le A_{n+1}(t)$
since its correlation measure is positive
and $\cos(xw)\le1$.
Multiplying the lower bound by $c_n(t)>0$
and subtracting it from the last upper bound gives
\begin{align*}
&k_nL_{n+1,t}(x)-c_n(t)L_{n,t}(x)\\
&\quad\le k_nA_{n+1}(t)
-4k_nA_{n+1}(t)\left(1-\frac{x^2}{48}\right)\\
&\quad=k_nA_{n+1}(t)\left(\frac{x^2}{12}-3\right).
\end{align*}
This proves \eqref{eq:sixbandresidual}.
The right side is nonpositive when $|x|\le6$,
giving \eqref{eq:sixbandsource}.
\end{proof}

\begin{corollary}[A global source estimate with a
controlled nonzero defect]
\label{cor:sixbandduhamel}
Put $U_n(s,x)=L_{n,\,1/2-s}(x)$,
$T=1/2$, and use
$c_n(s)=4k_nA_{n+1}(1/2-s)/A_n(1/2-s)$.
For the nonnegative source
\[
E_n(s,x)=[k_nU_{n+1}(s,x)-c_n(s)U_n(s,x)]_+
\]
we have, for all $s$ and $x$,
\begin{equation}
0\le E_n(s,x)\le
k_nA_{n+1}(1/2-s)
\left(\frac{x^2}{12}-3\right)_+ .
\end{equation}
Consequently its heat-smoothed central mass obeys
\begin{equation}\label{eq:sixbandduhamel}
\boxed{
\int_0^{1/2}(P_{1/2-s}E_n(s,\cdot))(0)\,ds
\le\frac{k_n}{3456}
\sup_{0\le t\le1/2}A_{n+1}(t).
}
\end{equation}
\end{corollary}
\begin{proof}
The positive-part inequality follows from
\eqref{eq:sixbandresidual}.
For $Z_h\sim N(0,h)$, $\mathbb E Z_h^4=3h^2$.
Since
$(x^2/12-3)_+\le x^4/432$
for every real $x$, Gaussian averaging yields
\[
\left(P_h(x^2/12-3)_+\right)(0)
\le h^2/144.
\]
Integrating over $0\le s\le1/2$ gives
$\int_0^{1/2}(1/2-s)^2\,ds/144=1/3456$,
which proves \eqref{eq:sixbandduhamel}.
\end{proof}

This is an all-order, entire-heat-slab
\emph{theta-specific estimate}, stronger in
quantifiers than the earlier certified $n=1$
rectangle. It is an explicit bound
on a nonnegative global defect, \emph{not}
a proof that the defect vanishes.
The central Duhamel criterion
(Theorem~\ref{thm:onepointcriterion}) needs
zero negative-part mass, or equivalently
a penalty bound uniform as its penalty
parameter tends to infinity.
The finite fixed-coefficient bound
\eqref{eq:sixbandduhamel} does not
supply that conclusion: at $|x|>6$
the polynomial envelope becomes positive.
It therefore does not prove RH.

The exact arithmetic checks are in
\path{certificates/xi_heat_six_band_exact.py};
the algebraic gates are formalized in
\path{XiKernel/SixBandDivisionFree.lean}.
The theta moment inputs are the already
validated two-resolution quadratures,
not new independently recalculated intervals.
The full infinite-kernel integral
identities and the semigroup estimate
are analytic rather than end-to-end Lean proofs.
No priority claim is made for the
curvature-to-correlation method;
Newman's earlier result is credited.
